% D2-2
% Auto-generated from markdown source
% origin: ai

\documentclass[12pt, a4paper]{article}
\usepackage[utf8]{inputenc}
\usepackage[T1]{fontenc}
\usepackage{amsmath, amssymb, amsfonts, amsthm}
\usepackage{textcomp}
\usepackage{geometry}
\usepackage{hyperref}
\usepackage{booktabs}
\usepackage{graphicx}
\usepackage{fancyhdr}
\usepackage{tcolorbox}
\usepackage{enumitem}

\geometry{margin=1in}
\pagestyle{fancy}
\fancyhf{}
\rhead{D2-2}
\lhead{\thepage}

\newtheorem{definition}{Definition}[section]
\newtheorem{theorem}{Theorem}[section]
\newtheorem{corollary}{Corollary}[section]
\newtheorem{proposition}{Proposition}[section]
\newtheorem{lemma}{Lemma}[section]
\newtheorem{conjecture}{Conjecture}[section]

\newcommand{\calI}{\mathcal{I}}
\newcommand{\calP}{\mathcal{P}}
\newcommand{\calE}{\mathcal{E}}
\newcommand{\calS}{\mathcal{S}}
\newcommand{\calH}{\mathcal{H}}
\newcommand{\calO}{\mathcal{O}}
\newcommand{\calF}{\mathcal{F}}
\newcommand{\calU}{\mathcal{U}}
\newcommand{\calV}{\mathcal{V}}
\newcommand{\calA}{\mathcal{A}}
\newcommand{\calG}{\mathcal{G}}
\newcommand{\calB}{\mathcal{B}}
\newcommand{\calC}{\mathcal{C}}
\newcommand{\calL}{\mathcal{L}}
\newcommand{\calM}{\mathcal{M}}
\newcommand{\calX}{\mathcal{X}}
\newcommand{\calY}{\mathcal{Y}}
\newcommand{\calW}{\mathcal{W}}
\newcommand{\calT}{\mathcal{T}}
\newcommand{\calZ}{\mathcal{Z}}
\newcommand{\calR}{\mathcal{R}}
\newcommand{\calQ}{\mathcal{Q}}
\newcommand{\calD}{\mathcal{D}}
\newcommand{\calN}{\mathcal{N}}
\newcommand{\calK}{\mathcal{K}}

\newcommand{\Tr}{\text{Tr}}
\newcommand{\E}{\mathbb{E}}
\newcommand{\diff}{\mathrm{d}}
\newcommand{\KL}{D_{\mathrm{KL}}}
\newcommand{\argmax}{\arg\max}
\newcommand{\argmin}{\arg\min}
\newcommand{\hbarF}{\hbar}

\newtcolorbox{thmbox}[1]{
  colback=blue!5!white, colframe=blue!75!black, fonttitle=\bfseries,
  title={#1}, breakable
}
\newtcolorbox{defbox}[1]{
  colback=green!5!white, colframe=green!75!black, fonttitle=\bfseries,
  title={#1}, breakable
}

\begin{document}

\title{D2-2}
\author{RALPH Loop}
\date{}
\maketitle

\begin{abstract}
This paper presents D2-2, formalizing the underlying principles using information-theoretic methods.

\textbf{Keywords}: D2-2, information theory, formalization.
\end{abstract}

\tableofcontents
\section{D2-2}




\#\#



\#\#

**** Agent  \(K_0\) \(\Delta K_t\)""

\[T_{emerge} = \inf\{t : I(K_t; y_{target}) > I(K_0; y_{target}) + \gamma\}\]

\textbf{*}*** \(t_c\) \(\sum_{i=0}^{t_c} \Delta K_i\) Agent """"


\#\#

\subsubsection{ 1}


****
\[\frac{dK_t}{dt} = -\lambda K_t + \alpha \cdot \nabla_K I(K; y_{target}) + \eta(t)\]


\begin{itemize}
\item \(-\lambda K_t\)****
\item \(\alpha \cdot \nabla_K I(K; y_{target})\)****
\item \(\eta(t)\)****
\end{itemize}


****\(\frac{dK_t}{dt} = 0 \implies \lambda K^* = \alpha \cdot \nabla_K I(K^*; y_{target})\)


\subsubsection{ 2}


****Cahn-Hilliard
\[\det\left(\frac{\partial^2 V_{eff}}{\partial K^2}\right) = 0\]

\(V_{eff}(K)\)
\[V_{eff}(K) = \frac{\lambda}{2} K^2 - \alpha I(K; y_{target}) + \frac{\sigma^2}{2} \log |\det(\nabla^2 K)|\]

****
\[t_c = \frac{1}{\lambda} \log\left(\frac{K^* - K_0}{K^* - K_c}\right)\]

\(K_t > K_c\) \textbf{*}***


\subsubsection{ 3D2-1/D2-2}


\textbf{D2-1  + D2-2  = }

\[\text{} \equiv |A| > |A|_c \land t > t_c\]

****
\begin{itemize}
\item D2-1Agent **** \(|A| = |A|_c\)
\item D2-2Agent **** \(t = t_c\)
\end{itemize}


****
\[P_{\text{emergence}}(|A|, t) = \Phi\left(\frac{|A| - |A|_c}{|A|_c}\right) \cdot \Phi\left(\frac{t - t_c}{t_c}\right)\]

\(\Phi\)


\subsubsection{ 4 D9-III}


\textbf{D9-III }
\begin{itemize}
\item D9-III  \(T_{train}\)
\item D2-2 \(T_{train} > t_c\) ****
\end{itemize}


\textbf{D9-III + D2-2 }
\[\text{} \equiv \mathbb{I}[|A| > |A|_c] \land \mathbb{I}[t > t_c] \land \mathbb{I}[I_{gen} > I_{cov}]\]

\textbf{** Agent }**


\#\#

\subsubsection{ 1Agent}

```python
\section{1.  Agent}

\section{2.  I(K_t; y_target)}

\section{3.  t_c}

```

\subsubsection{ 2RAG vs  RAG}

```python
\section{1.  RAGvs  RAG}

\section{2. RAG  t_c |A|_c}

```

\subsubsection{ 3}

```python
\section{1.}

\section{2.  t_c vs L}

```


\#\#

1. ****
2. \textbf{ D2-1 }
3. \textbf{ D9-III  \(T_{train}\) }


\#\#

\begin{itemize}
\item [x]  1-4
\item [ ]  1Agent
\item [ ]  2RAG vs  RAG
\item [ ]  3
\end{itemize}



* D2-2  D2-1 *
*2026-05-04 | r01 | *

\begin{thebibliography}{99}
\bibitem{ref1} Shannon, C.E. (1948). A Mathematical Theory of Communication.
\bibitem{ref2} Cover, T.M. (2006). Elements of Information Theory.
\end{thebibliography}

\end{document}